\documentclass[12pt,a4paper,english]{article}
\usepackage[utf8]{inputenc}
\usepackage[english]{babel}
\usepackage{graphicx}
\usepackage{epsfig}            % To allow figures
\usepackage{pstricks}          % To draw color pictures directly
\usepackage{fancyhdr}
\usepackage[verbose,a4paper,tmargin=30mm,bmargin=37.5mm,lmargin=30mm,rmargin=30mm]{geometry}
\usepackage{helvet}
\usepackage{mathptmx}
\usepackage[T1]{fontenc}
\usepackage{rotating}
\usepackage{amssymb}       % More special characters          
\usepackage{amsmath}       % More mathematical characters
\usepackage{natbib} % Needed e.g. for \citep
\usepackage{multirow}
\usepackage{verbatim}          % Use the verbatim styles
\definecolor{METblue}{cmyk}{0.85,0,0.2,0.2}
\usepackage{titlesec}
\titleformat{\section}
  [block]
  { \normalfont\sffamily\Large\bfseries\color{METblue} }
  {\makebox[2em][r]{\thesection}}
  {5mm}
  {\vspace{5mm}}

\titleformat{\subsection}[block]%
  {\normalfont\sffamily\bfseries}%
  {\makebox[2em][r]{\thesubsection}}%
  {5mm}
  {\vspace{3mm}}[]
\titleformat{\subsubsection}[block]%
  {\normalfont\sffamily\bfseries}%
  {\makebox[3em][r]{\thesubsubsection}}%
  {5mm}
  {\vspace{3mm}}[]

% Titlespacing syntax: 
\titlespacing*{\section}{-17mm}{5mm}{0mm}
\titlespacing*{\subsection}{-13.5mm}{5mm}{0mm}
\titlespacing*{\subsubsection}{-17.5mm}{5mm}{0mm}

%%% Formatting the table of contents
\usepackage{tocloft}
\renewcommand{\cfttoctitlefont}{\sffamily\bfseries\color{METblue}\Large}
\providecommand{\cftchapfont}{\sffamily\bfseries }
\renewcommand{\cftsecfont}{\sffamily\bfseries }
\renewcommand{\cftsubsecfont}{\sffamily }
\renewcommand{\cftsubsubsecfont}{\sffamily }
\providecommand{\cftsubsubsubsecfont}{\sffamily }
\renewcommand{\cftfigfont}{Figure }
\renewcommand{\cfttabfont}{Table }
\providecommand{\cftchappagefont}{\sffamily}
\renewcommand{\cftsecpagefont}{\sffamily\bfseries}
\renewcommand{\cftsubsecpagefont}{\sffamily}
\renewcommand{\cftsubsubsecpagefont}{\sffamily}
\providecommand{\cftsubsubsubsecpagefont}{\sffamily}

\renewcommand{\baselinestretch}{1.33}

% --------------------------------------

\begin{document}

\bibliographystyle{agufull04}
%\bibliographystyle{ams}
\thispagestyle{empty}  % Hide page numbers

\noindent
\begin{tabular}{@{} p{63mm} p{50mm} r}
\includegraphics*[]{met_rapport_logo_eng} % Automatically uses PDF or EPS in same directory depending on latex or pdflatex.
&
\fontsize{27.5pt}{33pt} \selectfont \bf \sffamily MET{\color{gray} report}
&
 \begin{minipage}[b]{28mm}
  \begin{flushright}
   \footnotesize \sffamily No. 13/2016 \\ ISSN 2387-4201 \\ Oceanography              % Report number and Category
  \end{flushright}
 \end{minipage}
\end{tabular}


%vfill
\vspace{0.5cm}

\begin{flushright}
{\fontsize{28pt}{43.2pt}\selectfont \bf \sffamily Comparison of simulations and observations in the Oslofjord}  \\        % Title
{\fontsize{14.0pt}{16.8pt}\selectfont \bf \sffamily FjordOs technical report No. 3}   % Subtitle

{\fontsize{12.5pt}{15pt}\selectfont \sffamily June 2016                                          % Subtitle
\\
\sffamily Karina Hjelmervik$^1$, Nils Melsom Kristensen$^2$, and Andre Staalstr\o m$^3$           % Author name(s)
}
\end{flushright}

%\vspace{25mm}
%\vspace{5mm}

\begin{figure}[!h]
\begin{flushleft}
\includegraphics*[height=5.5cm,angle=7]{Figurer/Driftere_ombord} \\
\end{flushleft}
\vspace{-2.5cm}
\begin{center}
\includegraphics*[height=6cm,angle=-4]{Python/salthold_96.png} \\
\end{center}
\vspace{-11cm}
\begin{flushright}
\includegraphics*[height=6.5cm,angle=-10]{Figurer/CTD} 
\end{flushright}
\vspace{3.5cm}
\begin{center}
\rput[bl](-7.5,0){\includegraphics[height=1.2cm]{Figurer/logo_met}} 
\rput[b](0,-0.2){\includegraphics[height=1.7cm]{Figurer/logo_hsn}} 
\rput[br](7.5,0){\includegraphics[height=1.2cm]{Figurer/logo_niva}} 
\end{center}
\end{figure}
\vspace{-1cm}
\noindent$^1$University College of Southeast Norway, $^2$Norwegian Meteorological Institute,\\
$^3$Norwegian Institute for Water Research 

%\newpage

%\thispagestyle{empty}  % Hide page numbers

\clearpage

\setlength{\unitlength}{1mm}  %Needed for picture environment

\begin{table}[!ht]

\begin{tabular}[c]{lr}
\vspace{5mm}
\includegraphics*{met_rapport_logo_eng} & \hspace{43mm}
{\fontsize{27.5pt}{33pt}\selectfont \bf \sffamily MET{\color{gray} report}}\\
\end{tabular}

\sffamily{
\begin{tabular}[t]{|p{110mm}|p{40mm}|} \hline
{\bf \sffamily Title}                  & {\bf \sffamily Date}               \\ 
Comparison of simulations and observations in the Oslofjord
                             & \today                   \\ \hline
{\bf \sffamily Section}                & {\bf \sffamily Report no.}         \\ 
                             &  13/2016                  \\ \hline
{\bf \sffamily Author(s)}                 & {\bf \sffamily Classification}     \\ 
Karina Hjelmervik, Nils Melsom Kristensen and Andre Staalstr\o m
                             & \begin{picture}(20,4)(-2,-1.0)
                               \put (0,0){\circle*{4}}
                               \put (7,0){\makebox(0,0){Free}}
                               \put (15,0){\circle{4}}
                               \put (27,0){\makebox(0,0){Restricted}}
                               \end{picture}
                               \\ \hline
{\bf \sffamily Client(s)}              & {\bf \sffamily Client's reference} \\ 
                                      &               \\ \hline
\end{tabular}

\begin{tabular}[t]{|p{154.3mm}|}
{\bf \sffamily Abstract}                                          \\
A three-day scientific cruise on board the research vessel (R/V) Trygve Braarud in the Oslofjord from the 21st to 23rd of September 2015 provided observations on hydrography and trajectories of drifters. The cruise was carried out as a part of the FjordOs project to investigate if targeted observations of hydrography and drifters would benefit the validation of the new Oslofjord model, the FjordOs-model. 
The cruise took place shortly after the storm "Petra" produced heavy rainfall in the areas surrounding the Oslofjord and caused an increase in river discharge. %Vertical CTD (Conductivity, Temperature and Depth) profiles at stations located along three different transects across the fjord reveal that the strong input of fresh water from Drammen river can be tracked halfway down the Oslofjord.

Provided is documentation of the observations and comparisons with simulations using the FjordOs-model. 
%Drifters were released in groups of three in two drop zones. 
%Four drifters were released along a line in a third drop zone, and in addition two individual drifters were released at selected spots. 
The drifter trajectories are compared with simulated trajectories using the open source trajectory-model OpenDrift.
%\\[50mm] % Add whitespace if necessary
\\ \hline
{\bf \sffamily Keywords}                                          \\ 
  ROMS, Oslofjord, CTD, drifter \\ 
\hline
\end{tabular}
}

\begin{tabular}[t]{cc}
%                             &                            \\
                             &                            \\
                             &                            \\
\line(1,0){70}               & \line(1,0){70}             \\ 
Disciplinary signature       & Responsible signature      \\
%Jan Erik Haugen             & \O{}ystein Hov             \\       % Add names if needed
\hspace{75mm}                & \hspace{75mm}              \\

\end{tabular}
\end{table}

\clearpage

\thispagestyle{fancy} % footer from fancyhdr package
\headheight=15pt
\renewcommand{\headrulewidth}{0pt}

%\clearpage
\section*{\hspace{17mm}Abstract}
A three-day cruise on board the research vessel (R/V) Trygve Braarud in the Oslofjord from the 21st to 23rd of September 2015 provided observations on hydrography and trajectories of drifters. The cruise was carried out as a part of the FjordOs project to investigate if targeted observations of hydrography and drifters would benefit the validation of the new Oslofjord model, the FjordOs-model.
The cruise took place shortly after the storm "Petra" produced heavy rainfall in the areas surrounding the Oslofjord and caused an increase in river discharge. 
%Vertical CTD (Conductivity, Temperature and Depth) profiles at stations located along three different transects across the fjord reveal that the strong input of fresh water from Drammen river can be tracked halfway down the Oslofjord.

Provided is documentation of the observations and comparisons with simulations using the FjordOs-model. 
%Drifters were released in groups of three in two drop zones. Four drifters were released along a line in a third drop zone, and in addition two individual drifters were released at selected spots. 
The drifter trajectories are compared with simulated trajectories using the open source trajectory-model OpenDrift.

\clearpage

\vfill

\fancyfoot{
% If abstract on separate page is not needed, move the following table to the page before
\begin{tabular}[b]{p{40mm}p{25mm}p{25mm}p{25mm}p{25mm}}
 \begin{minipage}[l]{40mm} \tiny \color{METblue} {\bf Norwegian Meteorological Institute}\\ Org.no 971274042\\ post@met.no\\ www.met.no / www.yr.no
 \end{minipage} & 
 \begin{minipage}[l]{25mm} \tiny \color{METblue} {\bf Oslo}\\ P.O. Box 43, Blindern\\ 0313 Oslo, Norway\\ T. +47 22 96 30 00
 \end{minipage} &
 \begin{minipage}[l]{25mm} \tiny \color{METblue} {\bf Bergen}\\ All\'egaten 70\\ 5007 Bergen, Norway\\ T. +47 55 23 66 00
 \end{minipage} & 
 \begin{minipage}[l]{25mm} \tiny \color{METblue} {\bf Troms\o}\\ P.O. Box 6314, Langnes\\ 9293 Troms\o, Norway\\ T. +47 77 62 13 00
 \end{minipage} & 
 \begin{minipage}[l]{25mm} \tiny \color{METblue} 
 \end{minipage}
\end{tabular}
}

\clearpage
\tableofcontents
\newpage 
\listoffigures 
\newpage
\listoftables

\clearpage
% Begin arabic page numbering (starts on 1)
%\pagenumbering{arabic}

\section{Introduction}
Provided is a documentation of a scientific cruise in the Oslofjord from the 21st to 23rd of September 2015. The cruise was a pilot study to investigate if targeted observations would benefit the validation of the new Oslofjord model \citep{roed:etal:2016}, and was carried out as part of the FjordOs project. FjordOs is a cooperation between MET Norway, University College of Southeast Norway (HSN), The Norwegian Institute for Water Research (NIVA), The Norwegian Coastal Administration (Kystverket), Exxonmobil, Norwegian Defence Research Establishment (FFI), Vestfold, Buskerud, and \O stfold county, and AGNES AB Milj\o konsulent.

Nine drifters were used during the cruise in addition to CTD measurements (Figure~\ref{fig:Utstyr}). The drifters were dropped in four different zones. CTD measurements were taken along three different transects across the fjord. %The ships log were used to estimate the current velocity by finding the difference between speed over ground and speed through water measured at 2.2 meters depth below the keel of the vessel. 
The drifter trajectories and the CTD measurements were compared with model results using the FjordOs model \citep{roed:etal:2016}. 
Drifters may also be used to estimate diffusivity. \cite{lacas:etal:2014} estimated diffusivities using a line of drifters released at the same latitude. The same approach applies with any ensemble, but requires more drifters than in this pilot study.


\begin{figure}[b]
\centerline{
%\includegraphics*[width=0.65\textwidth]{Figurer/Driftere_ombord}
%\includegraphics*[width=0.36\textwidth]{Figurer/CTD}
\includegraphics*[height=7.5cm]{Figurer/Driftere_ombord}
\includegraphics*[height=7.5cm]{Figurer/CTD}
}
\caption{\small
The nine drifters before they were released (left) and the CTD (right).}
\label{fig:Utstyr}
\end{figure}

\begin{figure}[tb]
\centerline{
\includegraphics*[width=\textwidth]{Figurer/Vannstand}}
\caption{\small
Observed water level at Oscarsborg in September 2015 together with the time of the drifter drops. Water level data obtained from www.sehavniva.no.}
\label{fig:Waterlevel}
\end{figure}

\begin{figure}[tb]
\centerline{
\includegraphics*[width=\textwidth]{Figurer/River}}
\caption{\small
Observed flow rate of the Drammen river measured in Mj\o ndalen September 2015. Raw data obtained from NVE.}% The data is not quality assured.}
\label{fig:River}
\end{figure}

\begin{figure}[tb]
\centerline{
\includegraphics*[width=\textwidth]{Figurer/Gullholmen_vind_sep2015}}
\caption{\small
The hourly mean (black) and maximum (red) wind speed at Gullholmen measurement station during the time of the cruise. Data obtained from www.yr.no.}
\label{fig:Wind}
\end{figure}


\subsection{Weather conditions}

The storm "Petra" hit Buskerud, Vestfold, Telemark, and Agder from Monday 14th to Friday 18th of September 2015. The storm was caused by a main low pressure in the English channel which moved northwards crossing the North Sea. The low pressure caused an increased water level in the Oslofjord from the 17th to 18th of September before the water level was normalized around the 20th of September (Figure~\ref{fig:Waterlevel}).

Due to the heavy rainfall that followed, rivers and streams were flooded and caused an increased fresh water flow into the fjord. There are two large rivers connected with the Oslofjord, one in Drammen and the other in Fredrikstad with a mean discharge of 317 and 729 m$^3$/s respectively \citep{milli:etal:2011}. The river in Drammen had a larger discharge than normal during the first part of September, and an increase from around 650 m$^3$/s to around 1550 m$^3$/s from the 15th to 18th of September (Figure~\ref{fig:River}). The volume flux slowly decreased through the rest of September. 

During the cruise, the weather was sunny with light or gentle breeze (Figure~\ref{fig:Wind}). The wind direction Monday afternoon was southerly, turning east. 

\subsection{Ocean model}
The FjordOs-model was set up for the time of the cruise. The FjordOs model is curviliniar, free-surface, and terrain-following model based on the Rutgers Regional Ocean Modeling System (ROMS) \citep{haidv:etal:2008,shche:mcwil:2003,shche:mcwil:2005,shche:mcwil:2009} adapted to the Oslofjord \citep{roed:etal:2016}. 

Current fields at 2 meters depth (Figure~\ref{fig:Current1} - \ref{fig:Current4}) reveal a strong outflow from the Drammens fjord during the cruise. When comparing this to the sea surface salinity (SSS) in Figure~\ref{fig:Salinity1}, we conclude that this strong current is indeed a result of the increased freshwater discharge from the Drammen river.

%The simulated water level compared to the observerd water level...

\clearpage

% ----------------------------------------------------------------------------
\section{Observations}


\subsection{Drifter trajectories}
The drifters used during the cruise (see Figure~\ref{fig:Utstyr}) are prototype low-cost units produced by a norwegian startup. They are made from a plastic drainpipe 1m long, and is equipped with a foam "donut" approximately 20 cm from the top. The purpose of the donut is to stabilize the drifter, and to decrease vertical motion due to small surface waves. The pipe is filled with water when it is dropped into the sea, and is  ballasted for the drifter to stay vertically. The ballast is tuned so that the donut is level with the water surface.

The Arduiono\footnote{https://www.arduino.cc/}-based control unit of the drifters, located at the top of the drifter, is equipped with a GPS-reciever and a GSM-transmitter. The refresh rate of the positions sent to our database was set to 0.0167 Hz (updates once every minute). Unfortunately, some of the drifters failed\footnote{Stopped sending data to our database. The developer of the drifters has now fixed this failure.} during the cruise.

The drifters were released in groups of three in two drop zones (Figure~\ref{fig:ZoneA} and \ref{fig:ZoneB}), and four drifters along a line in a third drop zone (Figure~\ref{fig:ZoneC}). In addition two individual drifters were released (Figure~\ref{fig:ZoneD}). For more details see Table~\ref{tab:Drifters}. 




\begin{table}[tb]
%\vspace{-1.5cm}
\caption{Drifters}
\label{tab:Drifters}
\begin{tabular}{|c|c|ccc|} 
        \hline
        Drop & Drifter & Drop pos.            & Last sent pos.          & Picked up                  \\
        ID  & No.     & and time [UTC]       & and time [UTC]          &                            \\ \hline
             &         &          \multicolumn{3}{c|}{\textit{Drop zone A - Filtvet}}                 \\ 
        A$_1$& 1       & 59.548130, 10.617290 & 59.612400, 10.633990    & 59.581355, 10.619952       \\
             &         & 12:35 21.09.15       & 01:08 22.09.15          & 15:55 22.09.15             \\ 
        A$_2$& 8       & 59.548328, 10.617360 & 59.586281, 10.616590    & Found later in Hurum                  \\
             &         & 12:26 21.09.15       & 09:42 22.09.15          &                            \\
        A$_3$& 9       & 59.548351, 10.617710 & 59.527760, 10.537500    & In last sent position      \\
             &         & 12:22 21.09.15       & 16:27 22.09.15          &                            \\ \hline 
             &         &          \multicolumn{3}{c|}{\textit{Drop zone B - R\o dtangen}}             \\
        B$_1$& 3       & 59.531219, 10.405220 & 59.507019, 10.426410    & 59.494862, 10.476880       \\
             &         & 13:35 21.09.15       & 15:05 21.09.15          & 17:54 21.09.15             \\
        B$_2$& 5       & 59.524269, 10.404090 & 59.510750, 10.428240    & In last sent position     \\
             &         & 13:39 21.09.15       & 09:09 22.09.15          &                            \\
        B$_3$& 6       & 59.524261, 10.404080 & 59.507210, 10.546380    & 59.508735, 10.433803       \\
             &         & 13:39 21.09.15       & 22:37:15 21.09.15       & 09:12 22.0915              \\
        B$_4$& 10      & 59.494839, 10.477530 & 59.519531, 10.410070    & In last sent position      \\
             &         & 17:53 21.09.15       & 09:14 22.0915           &                            \\ \hline 
             &         &          \multicolumn{3}{c|}{\textit{Drop zone C - Horten-Moss}}             \\
        C$_1$& 4       & 59.430672, 10.514930 & 59.434101, 10.510170    & 59.437445, 10.507031       \\
             &         & 10:18 22.09.15       & 12:31 22.09.15          & 14:52 22.09.15             \\
        C$_2$& 5       & 59.431519, 10.525370 & 59.428638, 10.511890    & Found later in Hurum             \\
             &         & 10:26 22.09.15       & 15:53 22.09.15          &                            \\
        C$_3$& 10      & 59.430859, 10.540190 & 59.439060, 10.496760    & In last sent position      \\
             &         & 10:46 22.09.15	      & 08:50 23.09.15          &                            \\
        C$_4$& 7       & 59.432251, 10.557540 & 59.428982, 10.555320    & 59.406315, 10.532240       \\
             &         & 11:01 22.09.15       & 11:24 22.09.15          & 14:30 22.09.15             \\ \hline 
             &         &          \multicolumn{3}{c|}{\textit{Single drops}}                          \\
        D$_1$& 9       & 59.528332, 10.593090 & 59.523029, 10.520360    & 59.524059 10.520618        \\
             &         & 16:46 22.09.15       & 12:05 23.09.15          & 12:40 23.09.15             \\
%        13   & 3+6+7   & 59.505730, 10.424240 & 59.512871,  10.432890   & 59.454636 10.424795        \\
%             &         & 17:56 22.09.15       & 21:27 22.09.15          &                            \\
        D$_2$& 6       & 59.505692, 10.424270 & 59.474941, 10.431430    & 59.454636 10.424795        \\
             &         & 17:56 22.09.15       & 06:00 23.09.15          &                            \\ 
\hline 
\end{tabular}
\end{table}
\clearpage


% ----------ZONE A------------------------------------------------------------
\subsubsection{Drop zone A: Filtvet}
%
\begin{figure}[ht]
\centerline{
\includegraphics*[width=0.6\textwidth]{Figurer/zoneA}}
\caption{\small
The path and velocity of the drifters droppped in zone A.% together with the SOG estimated from the ship log.
}
\label{fig:ZoneA}
\end{figure}
%
% \begin{figure}[ht]
% \centerline{
% \includegraphics*[width=0.495\textwidth]{Opendrift_simulations/dropzone_A_no1_fjordos_arome}
% \includegraphics*[width=0.495\textwidth]{Opendrift_simulations/dropzone_A_no1_norkyst_arome}
% }
% \caption{\small
% Simulated trajectories of drifter dropped in zone A using OpenDrift (as described in Section \ref{sect:trajmod}). Left panel is forced by the FjordOs-CL1 model, and the right panel is forced by the NorKyst-800m model. Both have wind drift from the atmospheric model Arome-2.5km, but no wave drift is applied. Particles in OpenDrift were released over a period of 8 hours.
% }
% \label{fig:opendriftA}
% \end{figure}
%
\begin{figure}[ht]
\centerline{
\includegraphics*[width=0.495\textwidth]{Figurer/Driftere_vann2}
\includegraphics*[width=0.495\textwidth]{Figurer/Driftere_vann}
}
\caption{\small
The drifters shortly after releasing in zone A (left) and zone B (right).
}
\label{fig:DriftereVann}
\end{figure}

Three drifters were dropped in zone A approximately at 12:30 local time (LT is UTC+1) on 21st of September 2015. The drifters were dropped within two meters of each other (Figure~\ref{fig:DriftereVann}, left).

At the time of deployment, the currents were weak due to high water (Figure~\ref{fig:Current2}). The drifters stayed in approximately the same positions during ebb tide. At low water around 19:00LT, the drifters started to drift northwards and continued for the next 12 hours. The drifters had a maximum speed of 0.5 m/s. % Similar speeds is not found in the simulated current fields (Figure~\ref{fig:Current2}-\ref{fig:Current3}). 
The next morning the drifters had parted a bit, and the drifter closest to the coastline had a weaker velocity than the ones further out, and changed direction first (Figure~\ref{fig:ZoneA}). 

Two of the drifters were picked up the next day. One drifter stranded and was lost. It was later found by some locals and returned to the group.

% ----------ZONE B------------------------------------------------------------
\subsubsection{Drop zone B: R\o dtangen}
\label{sect:ZoneB}
%
\begin{figure}[ht]
\centerline{
\includegraphics*[width=0.9\textwidth]{Figurer/zoneB}}
\caption{\small
The path and velocity of the drifters droppped in zone B.}
\label{fig:ZoneB}
\end{figure}
%
% \begin{figure}[ht]
% \centerline{
% \includegraphics*[width=0.495\textwidth]{Opendrift_simulations/dropzone_B_no7_fjordos_arome}
% \includegraphics*[width=0.495\textwidth]{Opendrift_simulations/dropzone_B_no7_norkyst_arome}
% }
% \caption{\small
% As for Figure~\ref{fig:opendriftA}, but for drifter in zone B.
% }
% \label{fig:opendriftB}
% \end{figure}
%
Three drifters were dropped in zone B approximately at 13:30LT on 21st of September 2015. The drifters were dropped close to each other (Figure~\ref{fig:DriftereVann}, right panel).

Since there was little or no wind at the time of the deployment, we can give a crude estimate of the currents by estimating the ship drift. This is done by comparing the ship speed and course according to the GPS (Speed Over Ground, SOG, and Course Over Ground, COG) with speed through the water and compass heading. The difference between these two vectors give a estimate of the currents, which revealed that the drifters were deployed close to the eastern edge of the outflow of the Drammensfjord. The drifters followed the outflow from the Drammens fjord southwards before they turned east. The speed of the drifters were approximately 0.5 m/s (Figure~\ref{fig:ZoneB}). % which is approximately the same level as in the simulated current fields (Figure~\ref{fig:Current1}). 

One of the drifters stopped sending, and was replaced by a working unit at approximately 18:00LT. The drifters continued further east towards the small island Tofteholmen lying on a south-north threshold. At 22:00LT the drifters turned westwards. The speed of the drifters increased as they approached the outlet of the Drammensfjord. % This is in accordance with the simulated current fields (Figure~\ref{fig:Current3}). 
At approximately 09:00LT on the 22nd of September, the drifters were picked up north of the position they were deployed.

% ----------ZONE C------------------------------------------------------------
\subsubsection{Drop zone C: Horten-Moss}
%
\begin{figure}[ht]
\centerline{
\includegraphics*[width=0.9\textwidth]{Figurer/zoneC}}
\caption{\small
The path and velocity of the drifters droppped in zone C.}
\label{fig:ZoneC}
\end{figure}
%
% \begin{figure}[ht]
% \centerline{
% \includegraphics*[width=0.495\textwidth]{Opendrift_simulations/dropzone_C_no10_fjordos_arome}
% \includegraphics*[width=0.495\textwidth]{Opendrift_simulations/dropzone_C_no10_norkyst_arome}
% }
% \caption{\small
% As for Figure~\ref{fig:opendriftA}, but for drifter in zone C.
% }
% \label{fig:opendriftC}
% \end{figure}
%
Four drifters were released at four positions in a cross-section of the fjord between Horten and Moss on September 22nd 2015.% According to the simulated current fields (Figure~\ref{fig:Current3}) they should have floated southwards with stronger speed the further east. 

When estimating the ship drift, using the same method as described in Section \ref{sect:ZoneB}, this indicated a flow towards south-southwest in the whole cross section at the time of deployment. The currents were weakest on the western side of the cross-section, towards Horten. At 11:30LT, in a position close to Horten, the ship drift revealed a weak current towards north close to Horten. All four drifters floated towards southwest, but turned towards northeast after a few hours. The drifter closest to Horten (no. 4), changed direction at approximately 11:30LT, and the two further east (no. 5 and 10) at approximately 14:40LT. The drifter furthest east (no. 7) stopped sending its position shortly after deployment, and was picked up to the southwest of its initial position 3.5 hours after deployment.

Only drifter no. 10 continued to send its position until it was picked up. The drifter stranded in a shallow area just east of the island Veal{\o}s north of Horten.

One drifter was lost during the cruise (no. 5), but was found by locals in Hurum and returned to the research group.

% ----------ZONE D------------------------------------------------------------
\subsubsection{Single drops}
%
\begin{figure}[ht]
\centerline{
\includegraphics*[width=0.9\textwidth]{Figurer/zoneD}}
\caption{\small
The path and velocity of the single drifters droppped in Breiangen.}
\label{fig:ZoneD}
\end{figure}
%
% \begin{figure}[ht]
% \centerline{
% \includegraphics*[width=0.495\textwidth]{Opendrift_simulations/dropzone_D_no12_fjordos_arome}
% \includegraphics*[width=0.495\textwidth]{Opendrift_simulations/dropzone_D_no12_norkyst_arome}
% }
% \caption{\small
% As for Figure~\ref{fig:opendriftA}, but for single drifter in drop no. 12.
% }
% \label{fig:opendriftD1}
% \end{figure}
% %
% \begin{figure}[ht]
% \centerline{
% \includegraphics*[width=0.495\textwidth]{Opendrift_simulations/dropzone_D_no13_fjordos_arome}
% \includegraphics*[width=0.495\textwidth]{Opendrift_simulations/dropzone_D_no13_norkyst_arome}
% }
% \caption{\small
% As for Figure~\ref{fig:opendriftA}, but for single drifter in drop no. 13.
% }
% \label{fig:opendriftD2}
% \end{figure}
%
Two single drifters were released in Breiangen (Figure~\ref{fig:ZoneD}).

One drifter (no. 9) was released in the northeastern part of Breiangen at 16:46LT on 22nd of September 2015. It drifted slowly towards the northeast, then south, and then north before it followed an westward current along the coast. %The simulated current fields reveal a system of unorganized eddies with weak velocities in this area at that time.

One drifter (no. 6) was released south in the northwestern part of Breiangen at 17:56LT on 22nd of September 2015. It first made a anti-clockwise turn with low drifting velocities. % an accordance with the simulated current fields (Figure~\ref{fig:Current3}). 
At 01:30LT the drifter started going southwards with the outflow from the Drammensfjord. The velocity of the drifter was almost 0.4m/s east of Lang{\o} ya. % which is also in accordance with the simulated current fields (Figure~\ref{fig:Current4}).   

% ----------------------------------------------------------------------------

%\subsection{CTD}
%The CTD measurements were performed using an ...
%
%Measurements were performed in three crossections in addtion to some extra positions which are applied on a regular basis over several years (Figure~?).
%
%To be continued ...
%
%\begin{figure}[ht]
%\centerline{
%\includegraphics*[width=\textwidth]{Figurer/Salt_snitt}}
%\caption{\small
%Salinity in the three crossections.}
%\label{fig:CTD_Salt}
%\end{figure}

\clearpage

% ----------------------------------------------------------------------------
\subsection{Hydrography}
\label{sect:hydro}
Episodes of intense precipitation leads to fierce water flow in the rivers, 
like it is observed in Drammen River from September 17 (Figure~\ref{fig:River}).
Such episodes are important for the water quality in the fjord. 
The most obvious effect is that the salinity in the surface layer of the fjord 
is reduced, but the river water also carry suspended particles, organic matter, 
nutrients as well as others substances like for instance contaminants.       

Vertical profiles with a CTD (Conductivity, Temperature and Depth)
from (R/V) Trygve Braarud was measured
at stations located along three different transects across the fjord. 
The CTD instrument was of the type Seabird SBE 911 plus.
The stations are listed in Table \ref{tab:CTD} and shown in the map in 
Figure~\ref{fig:CTD_Kart}. 

The profiles of salinity and temperature reveal the different water masses in 
the fjord separated by topography. Figure~\ref{fig:profiles} show profiles from 
four different stations, Fl1 inside the Dr{\o}bak Sill, Im2 in the Dr{\o}bak Sound, 
Pj2 located between Horten and Jel{\o}ya and Sm1 outside Slagen. 
The salinity at station Fl1 show that the water mass below approximately 30 m 
are different than in the rest of the fjord. 
The sill depth at Dr{\o}bak is about 20 m. 
By zooming in on the depth range 90-160 m, it can also be seen that the water 
masses below approximately 115 m in the Dr{\o}bak Sound are different than at 
stations further out in the fjord. The sill depth between the Dr{\o}bak Sond and 
Breiangen is about 115 m. 

\begin{table}[tb]
%\vspace{-1.5cm}
\caption{Hydrography stations}
\label{tab:CTD}
\centerline{
\begin{tabular}{|c|c|c|c|c|} \hline
Station & Date and time & Latitude & Longitude & Salinity 0-5 m \\ 
code    &  [UTC]        &          &           & [psu] \\ \hline
Fl1 & Sep 21 2015 08:19 & 59.754066 & 10.574734 & 19.09\\ \hline
Im2 & Sep 21 2015 09:18 & 59.754066 & 10.628217 & 21.93\\ \hline
D-1 & Sep 21 2015 11:29 &           &           & 17.35 \vspace{-0.2cm} \\ 
    & Sep 22 2015 15:27 & 59.754066 & 10.405000 & 15.62\\ \hline
Me1 & Sep 21 2015 11:58 &           &           & 14.71 \vspace{-0.2cm} \\  
    & Sep 23 2015 07:57 & 59.754066 & 10.360000 & 12.72\\ \hline
Mf1 & Sep 21 2015 12:12 &           &           & 14.66 \vspace{-0.2cm} \\ 
    & Sep 23 2015 08:11 & 59.754066 & 10.391667 & 15.05\\ \hline
Mg1 & Sep 21 2015 12:28 &           &           & 19.00 \vspace{-0.2cm} \\
    & Sep 21 2015 15:26 &           &           & 16.15 \vspace{-0.2cm} \\
    & Sep 22 2015 15:45 &           &           & 15.84 \vspace{-0.2cm} \\	
    & Sep 23 2015 08:26 & 59.754066 & 10.422500 & 17.44\\ \hline
Mh1 & Sep 21 2015 12:46 &           &           & 19.31 \vspace{-0.2cm} \\ 
    & Sep 23 2015 08:45 & 59.754066 & 10.460833 & 16.96\\ \hline
Mi1 & Sep 21 2015 13:02 &           &           & 20.43 \vspace{-0.2cm} \\ 
    & Sep 23 2015 09:02 & 59.754066 & 10.495833 & 16.86\\ \hline
Mj1 & Sep 21 2015 13:18 &           &           & 19.86 \vspace{-0.2cm} \\ 
    & Sep 23 2015 09:14 & 59.754066 & 10.518333 & 16.96\\ \hline
Mk1 & Sep 21 2015 13:32 &           &           & 21.18 \vspace{-0.2cm} \\ 
    & Sep 23 2015 09:30 & 59.754066 & 10.550000 & 19.83\\ \hline
Ml1 & Sep 21 2015 13:46 &           &           & 20.11 \vspace{-0.2cm} \\ 
    & Sep 23 2015 09:44 & 59.754066 & 10.581667 & 15.68\\ \hline
Mm1 & Sep 21 2015 14:02 &           &           & 19.15 \vspace{-0.2cm} \\ 
    & Sep 23 2015 10:00 & 59.754066 & 10.620000 & 14.74\\ \hline
Pi1 & Sep 22 2015 08:03 & 59.754066 & 10.508333 & 17.28\\ \hline
Pj1 & Sep 22 2015 08:25 & 59.431667 & 10.525000 & 15.91\\ \hline
Pj2 & Sep 22 2015 08:37 & 59.431667 & 10.540157 & 16.14\\ \hline
Pk2 & Sep 22 2015 08:53 & 59.432419 & 10.557180 & 17.30\\ \hline
Pl1 & Sep 22 2015 09:08 & 59.431667 & 10.573333 & 19.14\\ \hline
Sk1 & Sep 22 2015 10:24 & 59.311667 & 10.540000 & 21.74\\ \hline
Sl1 & Sep 22 2015 10:37 & 59.314167 & 10.570000 & 20.40\\ \hline
Sm1 & Sep 22 2015 10:55 & 59.316667 & 10.601667 & 17.77\\ \hline
Sn1 & Sep 22 2015 11:16 & 59.318333 & 10.628333 & 17.94\\ \hline
\end{tabular}}
\end{table}

\begin{figure}[ht]
\centerline{
\includegraphics*[width=\textwidth]{Figurer/Kart_Breiangen.png}}
\caption{\small
Positions of CTD stations (red dots) taken between September 21st and September 23rd 2015.
The colorbar indicate the water depth. The contour line for 120 m depth is 
drawn with a green line.}
\label{fig:CTD_Kart}
\end{figure}

\begin{figure}[ht]
\centerline{
\includegraphics*[width=\textwidth]{Figurer/Fig_profiles.png}}
\caption{\small
Profiles of salinity (left) and temperature (right) at four stations in the 
Oslofjord. The upper panels show the depth range from 0 to 165 m.
The lower panels show the depth range from 90 to 160 m, and station Fl1 is omitted.
}
\label{fig:profiles}
\end{figure}

The CTD instrument allows for other sensors to be installed on it. 
(R/V) Trygve Braarud has a Seabird SBE 911 plus CTD, 
with additional sensors for measuring oxygen concentration, turbidity, 
chlorophyll fluorescence and coloured dissolved organic matter (CDOM). 
All of these, except the oxygen sensor, are optical instruments. 
The turbidity is a measure of how much light is scattered in the water, 
which depends on particle concentration. The value of the unit FNU is 
approximately equal to a particle concertation of 1 mg/L, 
but the scattering can also be affected by other things than particles, 
for instance like water bobbles. Chlorophyll fluorescence is a measure of the amount 
of fluorescent light from algae, and is a proxy for the chlorophyll content 
in the water mass. However, the relation between chlorophyll fluorescence 
and actual chlorophyll concentration can vary considerably, 
for instance with depth. 
CDOM is also called yellow substance, and is the optically measurable 
component of the dissolved organic matter in water. 
There is usually a strong negative correlation between salinity and CDOM, 
since the rivers are the major source of organic matter to the coastal zone. 

During the cruise with (R/V) Trygve Braarud, three transects 
across the fjord was conducted with the CTD. Nine stations (Me1-Mm1) were taken 
across Breiangen from Lang{\o}ya in the west to Jel{\o}ya in the east. 
Five stations (Pi1-Pl1) were taken between Horten and Gullholmen. 
Four stations (Sk1-Sn1) were taken between Slagen and Edge{\o}ya. 
For each of the transect salinity, temperature, chlorophyll fluorescence, 
CDOM and oxygen saturation are shown (Figures \ref{fig:brei1} to \ref{fig:slagen}). 
Two features are striking, 
firstly the strong input of fresh water from Drammen river and secondly the algae bloom on the east side of the fjord.

On the first day of the cruise, the freshest water in Breiangen is found on the west 
side. On the east side the salinity is relatively high, and high fluorescence 
values are found at the station closest to Jel{\o}ya. 
Two days later the fresh water is more evenly distributed across Breiangen, 
except at stations Mj1 and Mk1. At these two stations the highest fluorescence 
in the surface layer values are found. It is very clear that the algae are not 
found in the freshest part of the river water, since the low salinity is less 
optimal for marine algae. The algae are nevertheless found close to the surface 
where more light is available, and where the water column are more stratified. 
It is clear that the CDOM and the salinity pattern are very similar, and that 
CDOM is a good indicator for river water. It can also be noted that high oxygen 
concentrations can be related to high concentrations of algae. 

Much of the same pattern as is seen in Breiangen, is found in the transect 
between Horten and Gullholmen. The water from the Drammen River is found in 
the western part, and an algae bloom is found at the station farthest to the east. 
At the transect between Slagen and Edge{\o}ya the picture is more diffuse, with both 
the freshest and the most fluorescence rich water to the east. 
The salinity in this area can just as well be influenced by water from Glomma, 
given the right atmospheric conditions, but we will here explain that the low 
salinity at this transect is caused by water from the Drammen river. 
First a mixing diagram can be shown (see Figure \ref{fig:CFdiagram}), where two water properties are plotted 
along each of the axes. Often temperature and salinity is chosen, 
but here we will use CDOM and chlorophyll fluorescence. 
If the low salinity at station Sm1 and Sn1 at 1.5 m is due to river water from 
Drammenselva, this water mass are a mixture of water with high CDOM and low 
fluorescence values and water with low CDOM and high fluorescence values. 
River water meets marine water with algae. 
The points D1 and D2 (Figures \ref{fig:brei1} and \ref{fig:brei2}) are chosen to 
represent the river water, 
and the points B1 and B2 are chosen to represent the marine water. 
From the mixture diagram (Figure~\ref{fig:CFdiagram}) it can be seen that it is 
plausible that the water masses at station Sm1 and Sn1 
(called S1 and S2 in Figure~\ref{fig:slagen}) indeed are a 
product of river water and marine water with algae. 
From the model results it is shown that it should be expected that the river 
water from Drammenselva is found at the east side of the Slagen transect
(see for instance Figure 27). 


\begin{figure}[ht]
\centerline{
\includegraphics*[width=\textwidth]{Figurer/Breiangen_21_09_2015_v2.png}}
\caption{\small
Salinity, temperature, coloured organic matter, chlorophyll fluorescense and oxygen 
saturation in the upper 10 meters of the Breiangen transect, taken on 
September 21st 2015. The point D1 represents the river water and the point B1
represents marine water with algae.}
\label{fig:brei1}
\end{figure}

\begin{figure}[ht]
\centerline{
\includegraphics*[width=\textwidth]{Figurer/Breiangen_23_09_2015_v2.png}}
\caption{\small
As for Figure \ref{fig:brei1}, but for September 23rd 2015.
The point D2 represents the river water and the point B2
represents marine water with algae.}
\label{fig:brei2}
\end{figure}

\begin{figure}[ht]
\centerline{
\includegraphics*[width=\textwidth]{Figurer/Horten_22_09_2015_v2.png}}
\caption{\small
As for Figure \ref{fig:brei1}, but for the Horten transect, taken on September 22nd 2015.}
\label{fig:horten}
\end{figure}

\begin{figure}[ht]
\centerline{
\includegraphics*[width=\textwidth]{Figurer/Slagen_22_09_2015_v2.png}}
\caption{\small
As for Figure \ref{fig:brei1}, but for the Slagen transect, taken on 
September 22nd 2015. The points S1 and S2 represent the freshest water 
masses in this transect.}
\label{fig:slagen}
\end{figure}

\begin{figure}[ht]
\centerline{
\includegraphics*[width=\textwidth]{Figurer/CF_diagram.png}}
\caption{\small
Mixing diagram for CDOM and chlorophyll fluorescence.
The points B1 and B2 represent the river water from Drammensfjorden.
The points B1 and B2 represent the marine water that contain algae in Breiangen.
The points S1 and S2 represent the freshest water masses in the transect 
across the fjord at the latitude of Slagen.}
\label{fig:CFdiagram}
\end{figure}

\clearpage

\section{Simulations}

\subsection{Trajectories}
\label{sect:trajmod}
To simulate the trajectories of the drifters, we have applied the open source trajectory-model OpenDrift. This is a trajectory model under development at MET Norway, and is described by its developers as "a software for modeling the trajectories and fate of objects or substances drifting in the ocean, or even in the atmosphere". It is distributed under a GPL v2.0 license, and is available on GitHub\footnote{https://github.com/knutfrode/opendrift}.

The OpenDrift model was forced with currents from the FjordOs model and with wind from the Arome-MetCoOp 2.5km (Arome2.5) atmospheric model (the same atmospheric model was used as forcing when running the FjordOs hindcast). We also performed OpenDrift-simulations using currents from the NorKyst-800m model together with wind from Arome2.5, which is the current operational setup for running trajectory simulations at MET Norway, to compare with the new FjordOs model. We can tune a number of parameters when running OpenDrift, e.g. random walk and the wind drift factor. Random walk was not used in our simulations, but variations in wind drift factor was applied. At each of the locations of the real drifters, we released six "virtual" drifters, from 2.5 hours before to 2.5 hours after the real drifters. This was done to see if the time of the release was an important factor. The resulting trajectories can be viewed in Figure~\ref{fig:opendrift1} - \ref{fig:opendrift4}.


\begin{figure}[ht]
\centerline{
\includegraphics*[width=0.495\textwidth]{Opendrift_simulations/LTR3/tokt_drifters_winddrift_0p0_radius_0_num_6_plusminus_2p5h_crop}
\includegraphics*[width=0.495\textwidth]{Opendrift_simulations/LTR3/tokt_drifters_winddrift_0p0_radius_0_num_6_plusminus_2p5h_norkyst_crop}
}
\caption{\small
To the left is the trajectories simulated using OpenDrift, forced by the FjordOs model and Arome-MetCoOp, and to the right forced by NorKyst-800m and Arome-MetCoOp. Wind drift factor was set to 0.0 in both simulations. Six drifters were released in the model at each of the locations and, from 2.5 hours before to 2.5 hours after, the times of the drops as described in Table \ref{tab:Drifters}.
}
\label{fig:opendrift1}
\end{figure}

\begin{figure}[ht]
\centerline{
\includegraphics*[width=0.495\textwidth]{Opendrift_simulations/LTR3/tokt_drifters_winddrift_0p1_radius_0_num_6_plusminus_2p5h_crop}
\includegraphics*[width=0.495\textwidth]{Opendrift_simulations/LTR3/tokt_drifters_winddrift_0p1_radius_0_num_6_plusminus_2p5h_norkyst_crop}
}
\caption{\small
Same as Figure~\ref{fig:opendrift1}, but using wind drift factor of 0.01.
}
\label{fig:opendrift2}
\end{figure}

\begin{figure}[ht]
\centerline{
\includegraphics*[width=0.495\textwidth]{Opendrift_simulations/LTR3/tokt_drifters_winddrift_0p2_radius_0_num_6_plusminus_2p5h_crop}
\includegraphics*[width=0.495\textwidth]{Opendrift_simulations/LTR3/tokt_drifters_winddrift_0p2_radius_0_num_6_plusminus_2p5h_norkyst_crop}
}
\caption{\small
Same as Figure~\ref{fig:opendrift1}, but using wind drift factor of 0.02.
}
\label{fig:opendrift3}
\end{figure}

\begin{figure}[ht]
\centerline{
\includegraphics*[width=0.495\textwidth]{Opendrift_simulations/LTR3/tokt_drifters_winddrift_0p5_radius_0_num_6_plusminus_2p5h_crop}
\includegraphics*[width=0.495\textwidth]{Opendrift_simulations/LTR3/tokt_drifters_winddrift_0p5_radius_0_num_6_plusminus_2p5h_norkyst_crop}
}
\caption{\small
Same as Figure~\ref{fig:opendrift1}, but using wind drift factor of 0.05.
}
\label{fig:opendrift4}
\end{figure}

\clearpage

\subsection{Hydrography}

In this chapter, we will compare how the salinity vary across the fjord in three 
transects, based on observations and model results.

CTD casts was taken at nine stations across Breiangen at two occasions, as described in Chapter \ref{sect:hydro}.
The first transect was taken on September 21st, and the second two days later 
(Figures \ref{fig:CTD_Mod_M1} and \ref{fig:CTD_Mod_M2}, upper panels).
At the first occasion the river water from the Drammen river is concentrated on the 
west side of the transect.
Two days later, the river water is more evenly distributed across Breiangen.
The model reflects this well
(Fig. \ref{fig:CTD_Mod_M1} and \ref{fig:CTD_Mod_M2}, lower panels).
The most striking difference between the observations and the model results,
are the too high salinity below the pycnocline in the model results (5-10 psu).
This could most probably be tracked to the initial and open boundary conditions.
This feature is seen in all the four transects shown here.

The transect from Horten to Moss was taken on September 22nd (Fig. \ref{fig:CTD_Mod_P}), and the most fresh surface water was found at the west side of this transect, just as in Breiangen.
This was the case both in the observations and in the model results.
The stratification was somewhat weaker in the middle of the fjord.
This can be seen by studying the salinity at 6 m depth where a minimum value in the 
middle of the fjord can be observed both in the observations and in the model results.

The transect between Slagen and Eld{\o}ya was taken on September 22nd (Fig. \ref{fig:CTD_Mod_S}).
In the observations the lowest surface salinity is found on the east side of the fjord,
in contrast to the other transects further north.
As shown earlier, this fresh water mass has its origin in the Drammen river.
This feature is not reflected in the model results at the same time.
This discrepancy could be due to wrong timing and that the most fresh 
water could be found at the east side of the transect, also in the model, 
at a different time.
Figure \ref{fig:Current1} and \ref{fig:Current2} show that a strong current is flowing from Drammensfjorden
and out of the fjord, that pass on the east side of the fjord outside Slagen.
This would explain the fresher water on this side in the CTD observations.
But at Sep 22 13:00, the current pattern is more irregular (see Figure {fig:Current4}). 

\begin{figure}[ht]
\centerline{
\includegraphics*[width=\textwidth]{Figurer/Transect_M_no1_salt.png}}
\caption{\small
Salinity in a transect across Breiangen. 
The upper panel shows observations with CTD taken
21/9-2016 14:00 LT. The lower panel shows model results from the same time.}
\label{fig:CTD_Mod_M1}
\end{figure}

\begin{figure}[ht]
\centerline{
\includegraphics*[width=\textwidth]{Figurer/Transect_M_no2_salt.png}}
\caption{\small
Salinity in a transect across Breiangen. 
The upper panel shows observations with CTD taken
23/9-2016 10:00 LT. The lower panel shows model results from the same time.}
\label{fig:CTD_Mod_M2}
\end{figure}

\begin{figure}[ht]
\centerline{
\includegraphics*[width=\textwidth]{Figurer/Transect_P_salt.png}}
\caption{\small
Salinity in a transect from Horten to Moss. 
The upper panel shows observations with CTD taken
22/9-2016 10:00 LT. The lower panel shows model results from the same time.}
\label{fig:CTD_Mod_P}
\end{figure}

\begin{figure}[ht]
\centerline{
\includegraphics*[width=\textwidth]{Figurer/Transect_S_salt.png}}
\caption{\small
Salinity in a transect from Slagen to Eld{\o}ya. 
The upper panel shows observations with CTD taken
22/9-2016 13:00 LT. The lower panel shows model results from the same time.}
\label{fig:CTD_Mod_S}
\end{figure}

\clearpage

% ----------------------------------------------------------------------------
\section{Discussion \& Conclusions}

The main objective of this scientific cruise was to assess whether we could make use of targeted observations of hydrography and drift-paths in the development and validation of the new FjordOs model covering the Oslofjord.

The main focus area for the observations was the Breiangen area, located approximately in the middle of the fjord in the north-south direction. This is one of the widest areas of the Oslofjord, except for the entrance itself, and is also very interesting because of the influence from the Drammensfjord, and Drammen river, which has a large freshwater flux that empties into the Breiangen area. Our cruise took place just after the heavy rain of the storm "Petra", and the impact of this increased freshwater flux was clearly visible in many of our measurements.
The average salinity in the upper 5 metres based on the CTD observations was 17.3 psu in Breiangen,
17.2 psu in the transect between Horten and Moss and 19.5 psu in the transect between 
Slagen and Eld{\o}ya.

Introducing wind drift in the simulated trajectories of the drifters is crucial to match and explain the observed drift. The best fit between simulations and observations was given by setting the wind drift factor to 1\% (the expected influence of the wind on the drifting buoys was expected to be between 0\% and 3\%). When we compare the trajectories of the drifters from our simulations to the ones that was observed, we can see that there is a very good agreement between the two in parts of the domain. Especially the drifters in Dropzone A and C (Figure \ref{fig:ZoneA} and \ref{fig:ZoneC}), and the single drops (in Figure \ref{fig:ZoneD}) showed a good comparison between observations and simulations. The drifters dropped near the mouth of the Drammensfjord (Dropzone B, Figure \ref{fig:ZoneB}) did not perform as well when compared to the observations. This could be connected to errors in the hydrography of the model, and that the water coming out of the Drammensfjord dit not spread out enough, but continued south across the Breiangen as a too narrow jet. The comparison of observed and modelled salinities in Figure \ref{fig:brei1} suggest that the freshwater layer on top of the water column is too shallow, which in turn could produce currents that are too surface intensified. It could also relate to wrong timing of the freshwater release from the Drammensfjord. The errors in hydrography are likely to be related to errors in the initial conditions of the model that comes from the operational NorKyst-800m model at MET Norway.


The FjordOs-model provided, in general, more realistic drifter trajectories than the operational NorKyst-800m model.


The number of observations in the Oslofjord is very sparse. When we consider the scales of the physical phenomenon, we try to resolve with the FjordOs-model, we can clearly see the need for targeted observations to validate and improve the model. It is important to have enough observations (both temporal and spatial) to give an estimate of the mean physical state of the fjord, both with regards to hydrography and currents. In addition, targeted observations in relation to e.g. special weather events have proven to be very useful to study specific situations, and how our model copes with these events.




\clearpage
\section*{\hspace{17mm}Acknowledgements}
The present research and development is mostly funded by the regional research fund Oslofjordfondet. We gratefully acknowledge the support. Lars Petter R\o ed at MET Norway has contributed with valuable discussions in the planning of the cruise and post processing of data. The authors would also like to thank the crew at the research vessel (R/V) Trygve Braarud. 

\clearpage
\section*{\hspace{17mm}Appendix}

\begin{figure}[ht]
\centerline{
\includegraphics*[width=1.0\textwidth]{Python/salthold_72}}
\caption{\small
Sea surface salinity from the FjordOs model at the start of the cruise. Note the freshwater coming out from the Drammensfjorden.
}
\label{fig:Salinity1}
\end{figure}

\begin{figure}[ht]
\centerline{
\includegraphics*[width=1.0\textwidth]{Python/salthold_96}}
\caption{\small
Sea surface salinity from the FjordOs model at approximately 24 hours into the cruise. Note the freshwater from the Drammensfjorden has spread out further compared to Figure~\ref{fig:Salinity1}.
}
\label{fig:Salinity2}
\end{figure}

\begin{figure}[ht]
\centerline{
\includegraphics*[trim=2.0cm 3cm 6.0cm 3.5cm,clip=true,height=7cm]{Python/stromfelt_70}
\includegraphics*[trim=3.7cm 3cm 1.3cm 3.5cm,clip=true,height=7cm]{Python/stromfelt_72}
}
\centerline{
\includegraphics*[trim=2.0cm 3cm 6.0cm 3.5cm,clip=true,height=7cm]{Python/stromfelt_74}
\includegraphics*[trim=3.7cm 3cm 1.3cm 3.5cm,clip=true,height=7cm]{Python/stromfelt_76}
}
\centerline{
\includegraphics*[trim=2.0cm 3cm 6.0cm 3.5cm,clip=true,height=7cm]{Python/stromfelt_78}
\includegraphics*[trim=3.7cm 3cm 1.3cm 3.5cm,clip=true,height=7cm]{Python/stromfelt_80}
}
\caption{\small
Simulated current fields at 2 meters depth every other hour during the cruise.}
\label{fig:Current1}
\end{figure}

\begin{figure}[ht]
\centerline{
\includegraphics*[trim=2.0cm 3cm 6.0cm 3.5cm,clip=true,height=7cm]{Python/stromfelt_82}
\includegraphics*[trim=3.7cm 3cm 1.3cm 3.5cm,clip=true,height=7cm]{Python/stromfelt_84}
}
\centerline{
\includegraphics*[trim=2.0cm 3cm 6.0cm 3.5cm,clip=true,height=7cm]{Python/stromfelt_86}
\includegraphics*[trim=3.7cm 3cm 1.3cm 3.5cm,clip=true,height=7cm]{Python/stromfelt_88}
}
\centerline{
\includegraphics*[trim=2.0cm 3cm 6.0cm 3.5cm,clip=true,height=7cm]{Python/stromfelt_90}
\includegraphics*[trim=3.7cm 3cm 1.3cm 3.5cm,clip=true,height=7cm]{Python/stromfelt_92}
}
\caption{\small 
Simulated current fields at 2 meters depth every other hour during the cruise.}
\label{fig:Current2}
\end{figure}

\begin{figure}[ht]
\centerline{
\includegraphics*[trim=2.0cm 3cm 6.0cm 3.5cm,clip=true,height=7cm]{Python/stromfelt_94}
\includegraphics*[trim=3.7cm 3cm 1.3cm 3.5cm,clip=true,height=7cm]{Python/stromfelt_96}
}
\centerline{
\includegraphics*[trim=2.0cm 3cm 6.0cm 3.5cm,clip=true,height=7cm]{Python/stromfelt_98}
\includegraphics*[trim=3.7cm 3cm 1.3cm 3.5cm,clip=true,height=7cm]{Python/stromfelt_100}
}
\centerline{
\includegraphics*[trim=2.0cm 3cm 6.0cm 3.5cm,clip=true,height=7cm]{Python/stromfelt_102}
\includegraphics*[trim=3.7cm 3cm 1.3cm 3.5cm,clip=true,height=7cm]{Python/stromfelt_104}
}
\caption{\small
Simulated current fields at 2 meters depth every other hour during the cruise.}
\label{fig:Current3}
\end{figure}

\begin{figure}[ht]
\centerline{
\includegraphics*[trim=2.0cm 3cm 6.0cm 3.5cm,clip=true,height=7cm]{Python/stromfelt_106}
\includegraphics*[trim=3.7cm 3cm 1.3cm 3.5cm,clip=true,height=7cm]{Python/stromfelt_108}
}
\centerline{
\includegraphics*[trim=2.0cm 3cm 6.0cm 3.5cm,clip=true,height=7cm]{Python/stromfelt_110}
\includegraphics*[trim=3.7cm 3cm 1.3cm 3.5cm,clip=true,height=7cm]{Python/stromfelt_112}
}
\centerline{
\includegraphics*[trim=2.0cm 3cm 6.0cm 3.5cm,clip=true,height=7cm]{Python/stromfelt_114}
\includegraphics*[trim=3.7cm 3cm 1.3cm 3.5cm,clip=true,height=7cm]{Python/stromfelt_116}
}
\caption{\small
Simulated current fields at 2 meters depth every other hour during the cruise.}
\label{fig:Current4}
\end{figure}



%\begin{table}[h]
%\vspace{-1.5cm}
%{\bf Table A1: ESD statistics.}\\
%\label{tab:1}
%\begin{tabular}{llll}
%\small OSLO - BLINDERN 18700 & 
%\small $T_m$ $R^2$= 71 -- 86 $T_x$ $R^2$= 69 -- 82 $T_n$ $R^2$= 49 -- 76 \\
%\small $\Delta T_m$ $q_{0.05}$= 3.21 $q_{0.50}$= 1.33 $q_{0.95}$= 1.75 \\  
%\small $\Delta T_x$ $q_{0.05}$= 3.09 $q_{0.50}$= 1.38 $q_{0.95}$= 2.16 \\  
%\small $\Delta T_n$ $q_{0.05}$= 3.15 $q_{0.50}$= 1.12 $q_{0.95}$= 1.01 \\
%\end{tabular}
%\end{table}

\clearpage
\pagebreak

\bibliography{referanse}
%\begin{thebibliography}{1}
%\end{thebibliography}

\clearpage
\pagebreak
 

\end{document}


